% 出席確認クイズ 第09回　可視化と探索的分析
%   attendance/attendance09.html から生成（解説は本ガイド用に加筆）。

\quizq{1}{時系列の推移を示すのに一般的に適したグラフはどれでしょうか?}%
  {ヒストグラム}%
  {円グラフ}%
  {散布図}%
  {\correct{折れ線グラフ}}%
  {正解は (d)。時間に沿った変化を示すには折れ線グラフが適しています。(a) ヒストグラムは分布、(b) 円グラフは構成比、(c) 散布図は2変数の関係を見るためのグラフです。グラフの種類は「何を伝えたいか」で選びます。}

\quizq{2}{2つの量的変数の関係を見るのに適したグラフはどれでしょうか?}%
  {棒グラフ}%
  {帯グラフ}%
  {\correct{散布図}}%
  {円グラフ}%
  {正解は (c)。散布図は、2つの量的変数の関係（相関の有無や傾向、外れ値）を視覚的に確認するのに適しています。(a) 棒グラフは量の比較、(b) 帯グラフと (d) 円グラフは構成比を示すためのものです。}

\quizq{3}{誤解を招く可視化の例として適切なものはどれでしょうか?}%
  {凡例を付ける}%
  {軸にラベルを付ける}%
  {単位を明記する}%
  {\correct{縦軸の原点を0にせず、差を誇張して見せる}}%
  {正解は (d)。縦軸の原点を0以外にした棒グラフは、わずかな差が大きな差のように見え、読み手を誤解させます。(a)(b)(c) の凡例・軸ラベル・単位の明記は、誤解を防ぐための良い実践です。ほかにも3D表示や二重軸、都合のよい期間の切り出しが誤解を招く例です。}

\quizq{4}{「探索的データ分析(EDA)」の説明として適切なものはどれでしょうか?}%
  {分析を外部に丸投げする}%
  {結論を先に決めてから、それに合う図だけを作る}%
  {\correct{仮説を固める前に、データを様々な角度から眺めて特徴を把握する}}%
  {データを見ずに報告書を書く}%
  {正解は (c)。探索的データ分析（EDA）は、本格的な分析や仮説の確定の前に、データの分布・傾向・異常値などを様々な角度から確認する作業です。(b) は結論を先に決めて都合のよい図だけを作る、確証バイアス（第11回）の典型例です。}

\quizq{5}{生成AIにグラフを作らせたときの人の役割として適切なものはどれでしょうか?}%
  {見た目が派手なものを選ぶ}%
  {グラフの種類は問わない}%
  {AIが作ったので確認は不要}%
  {\correct{目的に合っているか、誤解を招かないかを判断する}}%
  {正解は (d)。AIはグラフを素早く作れますが、目的に合っているか、誤解を招かないかの判断は人の役割です。(a) 派手さは判断基準ではなく、(c) AIが作ったグラフでも軸や単位の誤りはあります。AIに作らせ、人が解釈する、という役割分担が本講義の基本です。}
